\documentclass{ltc-book}

% ============================================================================
% 页面设置
% ============================================================================
\pageSetup{
  columns=20,
  lines=20,
  page-split-enabled=false,
}

% 内容设置 - 使用 font-name 指定字体（避免被自动字体检测覆盖）
\contentSetup{
  font-size=14pt,
  font-name=TW-Kai,
}

% 启用脚注功能（段末注模式）
\脚注设置{number-style=lujiao, indent=1em, spacing=0.5em}

% ============================================================================
% 页眉：非首页显示竖排书名
% ============================================================================
\newcount\firstpagenum
\firstpagenum=915

\AddToShipoutPictureFG{%
  \ifnum\value{page}>\firstpagenum
    \AtPageLowerLeft{%
      \put(\LenToUnit{141mm},\LenToUnit{192mm}){%
        \parbox[t]{2em}{%
          \centering\fontsize{9pt}{13pt}\selectfont
          史\\記\\卷\\十\\六%
        }%
      }%
    }%
  \fi
}

\begin{document}

\title{史記卷十六}
\chapter{秦楚之際月表第四}
% 设置起始页码（Lua 引擎 + LaTeX 计数器）
\设置页码{915}
\setcounter{page}{915}

\begin{正文}
史記卷十六

秦楚之際月表第四
\脚注{索隱張晏曰：「時天下未定，參錯變易，不可以年記，故列其月。」今案：秦楚之際，擾攘僭篡，運數又促，故以月紀事名表也。}

\输出脚注

太史公讀秦楚之際，曰：初作難，發於陳涉；虐戾滅秦，自項氏；撥亂誅暴，平定海內，卒踐帝祚，成於漢家。五年之間，號令三嬗。\脚注{集解音善。○索隱古「禪」字，音巿戰反。三嬗，謂陳涉、項氏、漢高祖也。}自生民以來，未始有受命若斯之亟\脚注{索隱音己力反。亟訓急也。}也。

\输出脚注

昔虞、夏之興，積善累功數十年，德洽百姓，攝行政事，考之於天，\脚注{集解韋昭曰：「謂舜受禪，在璇璣玉衡以齊七政。」}然後在位。湯、武之王，乃由契、后稷脩仁行義十餘世，不期而會孟津八百諸侯，猶以為未可，其後乃放弒。\脚注{索隱後乃放殺。殺音弒，謂湯放桀，武王討紂也。}秦起襄公，章於文、繆，獻、孝之後，稍以蠶食六國，百有餘載，至始皇乃能並冠帶之倫。以德若彼，\脚注{索隱即契、后稷及秦襄公、文公、穆公也。}用力如此，\脚注{索隱謂湯、武及始皇。}蓋一統若斯之難也。

\输出脚注

秦既稱帝，患兵革不休，以有諸侯也，於是無尺土之封，墮壞名城，銷鋒鏑，\脚注{集解徐廣曰：「一作『鍉』。」○索隱鏑音的。注「鍉」字亦音的。案：秦銷鋒鏑，作金人十二，以弱天下之兵也。}鉏豪桀，維萬世\脚注{索隱維訓度，謂計度令萬代安也。}之安。然王跡之興，起於閭巷，合從討伐，軼於三代，鄉秦之禁，適足以資賢者\脚注{索隱鄉秦之禁適足資賢者。鄉音向，許亮反。謂秦前時之禁兵景不封樹諸侯，適足以資後之賢者，即高帝也。言驅除患難耳。}為驅除難耳。故憤發其所為天下雄，\脚注{索隱指漢高祖。}安在無土不王。\脚注{集解白虎通曰：「聖人無土不王，使舜不遭堯，當如夫子老於闕里也。」}此乃傳之所謂大聖乎？\脚注{索隱言高祖起布衣，卒傳之天位，實所謂大聖。}豈非天哉，豈非天哉！非大聖孰能當此受命而帝者乎？

\输出脚注

\end{正文}
\end{document}
